% Generated by scripts/assemble.py; corrections belong in source/corrections.json.
\examhead{65}{西南大学 2025 年数学分析试卷}{65}
% p066-b12; source PDF page 66
\par\Needspace{5\baselineskip}
\blockanchor{p066-b12}
\majorquestion{一}计算题. 每题 10 分, 共 30 分.\par
% p066-b13; source PDF page 66
\blockanchor{p066-b13}
\question{1}$\lim_{n\to \infty}\left(\frac{1}{\ln(n + 1) - \ln n} -n\right).$\par
% p066-b14; source PDF page 66
\blockanchor{p066-b14}
\question{2}$\lim_{x\to0}\left(\frac{\cos x}{\cos 2x}\right)^{\frac{1}{\ln(1+x^{2})}}$\par
% p066-b15; source PDF page 66
\blockanchor{p066-b15}
\question{3}$\int_{-\pi}^{\pi}(\sin2x)(\cos(x^{2})+e^{x})\mathrm{d}x.$\par
% p066-b16; source PDF page 66
\par\Needspace{5\baselineskip}
\blockanchor{p066-b16}
\majorquestion{二}计算题. 每题 10 分, 共 30 分.\par
% p066-b17; source PDF page 66
\blockanchor{p066-b17}
\question{1}设 $F$ 有二阶连续偏导数, 且 $z = f(x, y)$ 是由 $F(x, x + y, x + y + z) = 0$ 所确定的隐函数, 求 $\frac{\partial^2 z}{\partial x^2}$.\par
% p066-b18; source PDF page 66
\par\Needspace{6\baselineskip}
\blockanchor{p066-b18}
\question{2}计算\par
\nopagebreak[4]
% p066-b19; source PDF page 66
\blockanchor{p066-b19}
\displayformula{\iiint_ {V} \left(x ^ {2} + y ^ {2} + z ^ {\frac {1}{2}}\right) \mathrm{d} x \mathrm{d} y \mathrm{d} z.}
% p066-b20; source PDF page 66
\blockanchor{p066-b20}
其中 $V$ 是由 $\left\{ \begin{array}{ll}z = y & \text{绕} z\text{轴旋转一周的曲面与} z = 1\text{所围成的区域}.\\ z = 0 & \end{array} \right.$\par
% p066-b21; source PDF page 66
\blockanchor{p066-b21}
\question{3}求 $f(x, y) = x^2 y (4 - x - y)$ 在由 $x$ 轴，$y$ 轴与 $x + y = 6$ 所围区域 $D$ 中的最大值与最小值.\par
% p066-b22; source PDF page 66
\par\Needspace{5\baselineskip}
\blockanchor{p066-b22}
\majorquestion{三}(15 分) 设 $f(x)$ 在 $(- \infty, a)$ 上连续, 且 $\lim_{x \to -\infty} f(x)$ 与 $\lim_{x \to a^{-}} f(x)$ 存在. 证明: $f(x)$ 在 $(- \infty, a)$ 上一致连续.\par
% p066-b23; source PDF page 66
\par\Needspace{5\baselineskip}
\blockanchor{p066-b23}
\majorquestion{四}(15 分) 设 $f(x)$ 在 $[a, b]$ 上连续, 在 $(a, b)$ 上可导, 且存在 $\xi \in (a, b)$, 满足 $f'(\xi) = 0$. 证明: 存在 $\eta \in (a, b)$, 使得 $f'(\eta) - f(\eta) + f(a) = 0$.\par
% p066-b24; source PDF page 66
\par\Needspace{5\baselineskip}
\par\Needspace{6\baselineskip}
\blockanchor{p066-b24}
\majorquestion{五}(15 分) 设 $f(x)$ 在 $[a, b]$ 上连续, 且对任意的 $x \in (a, b)$, 有 $f(x) > f(a) = f(b)$. 证明: 存在 $\xi \in (a, b)$, 使得\par
\nopagebreak[4]
% p066-b25; source PDF page 66
\blockanchor{p066-b25}
\displayformula{\frac {\int_ {a} ^ {\xi} f (x) \mathrm{d} x}{\xi - a} = \frac {\int_ {\xi} ^ {b} f (x) \mathrm{d} x}{b - \xi}.}
% p067-b00; source PDF page 67
\par\Needspace{5\baselineskip}
\blockanchor{p067-b00}
\majorquestion{六}(15 分) 设 $f(x) = \sum_{n=2}^{\infty} \frac{1}{n^x}$, $x \in (1, +\infty)$, 常数 $a > 1$. 判断 $\int_{a}^{+\infty} f(x) \, \mathrm{d}x$ 的敛散性并证明.\par
% p067-b01; source PDF page 67
\par\Needspace{5\baselineskip}
\par\Needspace{6\baselineskip}
\blockanchor{p067-b01}
\majorquestion{七}(15 分) 设 $f(x)$ 在 $[a, b]$ 上有连续的导函数, 记 $J = \frac{1}{b - a} \int_{a}^{b} f(x) \, \mathrm{d}x$. 证明:\par
\nopagebreak[4]
% p067-b02; source PDF page 67
\blockanchor{p067-b02}
\displayformula{\frac {1}{b - a} \int_ {a} ^ {b} | f (x) - J | \mathrm{d} x \leq \int_ {a} ^ {b} | f ^ {\prime} (x) | \mathrm{d} x.}
% p067-b03; source PDF page 67
\par\Needspace{5\baselineskip}
\blockanchor{p067-b03}
\majorquestion{八}(15 分) 证明: 含参量积分 $\int_{0}^{+\infty} \frac{y \cos xy}{1 + y^2} \, \mathrm{d}y$ 在 $(0, +\infty)$ 上不一致收敛.\par
